\documentclass[conference]{IEEEtran}
\IEEEoverridecommandlockouts
% ICDE 2027 Industry & Applications track (single-blind; 12 pages + unlimited refs;
% no appendices; AI-generated-content acknowledgement required).

\usepackage{cite}
\usepackage{amsmath,amssymb,amsfonts}
\usepackage{graphicx}
\usepackage{booktabs}
\usepackage{xcolor}
\usepackage[hidelinks]{hyperref}
\usepackage{textcomp}

\newcommand{\todo}[1]{\textcolor{red}{[TODO: #1]}}
\newcommand{\arcadedb}{ArcadeDB}

\begin{document}

\title{ArcadeDB: One ACID Engine for Documents, Graphs, Vectors, and Time Series\\
{\large\todo{working title; alternatives in .notes strategy doc}}}

\author{
\IEEEauthorblockN{Taewoon Kim}
\IEEEauthorblockA{HumemAI\\
taewoon@humem.ai}
\and
\IEEEauthorblockN{Roberto Franchini}
\IEEEauthorblockA{Arcade Data\\
r.franchini@arcadedata.com}
\and
\IEEEauthorblockN{Luca Garulli}
\IEEEauthorblockA{Arcade Data\\
l.garulli@arcadedata.com}
}

\maketitle

\begin{abstract}
\todo{Written last. Spine: multi-model unification as a real, adopted system,
one page/WAL/MVCC-transaction pipeline under documents, graphs, key-value,
time series, and vectors; every index type (LSM, full-text, geo, hash,
dense/sparse vector) commits inside the same transaction; Raft replication ships
model-agnostic WAL page diffs; one jar runs embedded in-process or as a server
speaking six wire protocols. Deep dive: what adding the vector model to a unified
substrate takes, including the deliberate transactional-source-of-truth /
eventually-consistent-ANN-graph tradeoff. Evaluation vs specialist engines +
lessons from a decade of multi-model engineering (OrientDB lineage).}
\end{abstract}

\begin{IEEEkeywords}
multi-model databases, ACID transactions, vector search, graph databases,
embedded databases
\end{IEEEkeywords}

\section{Introduction}\label{sec:intro}
\todo{Vector-era motivation: apps hold documents + relationships + embeddings of
the same entities; composed specialist stacks pay consistency/ops glue. Thesis:
unification at the storage layer, not the API gateway. Claims 1--4 from the
audit. Explicitly claim integration, not new algorithms.}

\section{Background and Design Goals}\label{sec:background}
\todo{PROSE PASS NEEDED. Raw co-author input below (Garulli, 2026-07-08).}

ArcadeDB descends from OrientDB, and its design goals are best understood as
deliberate corrections of three OrientDB limitations its authors lived with:
\begin{itemize}
  \item \textbf{Durability.} OrientDB's write-ahead log was asynchronous,
  single-threaded, and (in the authors' assessment) unreliable. ArcadeDB's
  journal is synchronous, parallel, and crash-proof. It follows the point-of-no-return
  discipline detailed in \S\ref{sec:architecture}.
  \item \textbf{Edge storage.} OrientDB managed edges with a tree-bag
  structure; ArcadeDB replaces it with a LIFO linked list that scales
  substantially better at high edge counts.
  \item \textbf{High availability.} The OrientDB multi-master model ``simply
  doesn't work'' (authors' words); ArcadeDB abandons it for a Raft
  leader/follower architecture (\S\ref{sec:ha}).
\end{itemize}
\todo{Design goals framing; what industry usage demanded. Third-party academic
adoption: DatAasee~\cite{himpe2024dataasee} (metadata catalog for research
libraries) and AptBacterialDB~\cite{aptbacterialdb} (ArcadeDB/openCypher graph
layer for an antibacterial-aptamer database), both independent multi-model
graph uses. Christian's "why we chose it" paragraph goes here or in Lessons.}

\section{Architecture: One Substrate, Five Models}\label{sec:architecture}

\arcadedb{} is organized around a single storage abstraction rather than a
federation of per-model engines. Everything that persists, whether it holds
documents, graph adjacency, key-value pairs, time-series points, or vector
index segments, is a \emph{paginated component}: a file of fixed-size pages
managed by one page manager, modified through one transaction context, and
made durable by one write-ahead log. The data models differ in the record
types and index structures layered on top; they do not differ in how bytes
reach disk, how changes become durable, or how replicas stay consistent. This
section walks that pipeline bottom-up and states precisely which guarantees
each model inherits from it.

\subsection{Pages, Transactions, and the Log}\label{sec:pipeline}

All reads and writes go through page-level copy-on-write. A transaction that
modifies any record obtains private copies of the affected pages from the page
manager and accumulates its changes there; concurrent readers continue to see
the unmodified originals. Commit is optimistic: at first phase, the engine
validates that the pages read by the transaction have not been changed by a
concurrent committer (a per-page version check), then serializes every dirty
page, \emph{including index pages}, into a single WAL entry. There is no
separate index log and no index-specific recovery path: a secondary-index
update, a graph edge insertion, and a document field change that happen in the
same transaction land in the same log record and either all survive a crash or
none do. Recovery replays the WAL forward from the last checkpoint; we measure
recovery time and post-crash integrity under \texttt{kill -9} in
\S\ref{sec:eval}.

Durability is configurable per deployment in one dimension only: when the WAL
buffer reaches the operating system (asynchronously on a background flusher by
default, or \texttt{fsync} per commit under the strictest setting). We flag
this knob here because it moves throughput by more than an order of magnitude
and is therefore the first thing our evaluation pins when comparing against
engines with different default contracts (\S\ref{sec:eval}).

\subsection{One Index Family Tree}\label{sec:indexes}

Table~\ref{tab:inventory} inventories the index and storage structures. The
workhorse is an LSM tree~\cite{oneil1996lsm} implemented as a paginated
component: mutable in-memory pages absorb writes, immutable compacted segments
serve reads, and compaction is an ordinary paginated write path. The full-text
index is a thin analyzer layer over the same LSM tree (tokens as keys), and
the geospatial index likewise reduces to LSM entries over space-filling-curve
keys. The practical consequence of this reuse is uniform transactional
behavior: because every index is a paginated component, the commit path of
\S\ref{sec:pipeline} covers all of them with no per-index machinery. The two
vector index families are the newest members of the tree and the most
demanding test of the claim that the substrate generalizes; they get their own
section (\S\ref{sec:vector}).

\subsection{Graph Storage, Stated Honestly}\label{sec:graph}

Vertices and edges are records in buckets, and adjacency is maintained as
RID-linked lists attached to each vertex, a design inherited from
OrientDB~\cite{ritter2021orientdb} and refined (the tree-bag edge container
was replaced by a LIFO linked list that behaves better at high fan-out,
\S\ref{sec:background}). Edges without properties are stored as
\emph{light edges}: pure RID pairs in the adjacency structure with no edge
record at all, which roughly halves the storage and I/O cost of plain
connections. Traversal is pointer-chasing over the record store through the
same page cache as every other access. We do not claim parity with
purpose-built native graph storage on every workload; we claim, and measure in
\S\ref{sec:eval}, that RID-linked adjacency over a shared substrate delivers
competitive point-traversal latency while inheriting transactions, recovery,
and replication that a bolted-on graph layer would have to re-implement.

\subsection{Replication Is Model-Agnostic}\label{sec:replication-arch}

Because every model's changes reduce to WAL'd page deltas, replication does
not know what a document, an edge, or a vector is. The high-availability
module (\S\ref{sec:ha}) ships transaction page deltas through a Raft
log~\cite{ongaro2014raft}; followers apply them through the same recovery
code path used after a crash. Adding a new data model to \arcadedb{}, as
\S\ref{sec:vector} recounts for vectors, required \emph{no} replication work:
a model whose writes are paginated is replicated by construction. The one
deliberate exception, derived search structures that are rebuilt rather than
replicated, is discussed in \S\ref{sec:vector}.

\section{Case Study: Adding the Vector Model}\label{sec:vector}

Vector search is the sharpest recent test of the unification thesis: it
arrived as a workload after the substrate was fixed, it is dominated by
specialist engines with no transactional obligations, and it stresses exactly
the properties (index build cost, memory, approximate results) that a shared
substrate might be suspected of handling poorly. This section describes what
adding dense and sparse vector search to \arcadedb{} actually took, in the
spirit of integration case studies such as
SingleStore-V~\cite{chen2024singlestorev} and Cosmos DB's DiskANN
integration~\cite{upreti2025cosmosvector}, and states the one place where we
deliberately relaxed the substrate's guarantees.

\subsection{Dense Vectors: Transactional Records, Derived Graph}

Dense embeddings are ordinary record properties: they live in buckets, are
WAL'd, versioned, and replicated like any other field. Approximate search is
served by a graph index built with the JVector
library~\cite{jvector2024repo}, which combines an HNSW-style layer
hierarchy~\cite{malkov2020hnsw} with Vamana-style graphs within each
layer~\cite{subramanya2019diskann}. Here we made the one deliberate exception
to full unification: the ANN graph itself is an asynchronously maintained,
WAL-disabled structure. Vector \emph{records} are the transactional source of
truth; the search graph is derived state, rebuilt from records if lost, and
followers rebuild rather than receive it. The alternative, WAL-logging every
neighbor-list mutation of an incremental graph build, would have coupled
commit latency of unrelated transactions to ANN maintenance. The consequence
is a bounded staleness window between a committed insert and its
visibility in approximate search (newly inserted vectors are meanwhile served
exactly from a delta buffer), which we consider the correct trade for an
operational engine and state openly rather than claim a fully transactional
ANN index. One design decision this enabled: because vectors-as-records are
the ground truth, index parameters can be changed and the graph rebuilt
without touching data, which we used in \S\ref{sec:eval} to match graph
degrees fairly across engines.

\subsection{Sparse Vectors: An Inverted Index That Commits}

Learned sparse retrieval (SPLADE-style~\cite{formal2021splade}) maps text to
high-dimensional sparse weight vectors and is conventionally served by
standalone inverted-index engines with block-max
pruning~\cite{broder2003wand,ding2011blockmax,tonellotto2018efficient}. To our
knowledge \arcadedb{} is the first general-purpose transactional DBMS to ship
this workload natively: \texttt{LSM\_SPARSE\_VECTOR} stores posting lists as
LSM segments (the same paginated component family as every other index), with
RID-delta compression, per-block weight quantization to int8, block maxima
for WAND-style pruning at query time, and size-tiered compaction as the
settle step. Because postings ride the LSM commit path, sparse index updates
are transactional with document writes, a property none of the standalone
engines offer. The quantization choice is a measured trade: int8 posting
weights cost about half a point of recall@10 against exact scoring
(\S\ref{sec:eval}) and can be disabled per index (fp32 segments) where exact
weights matter more than space.

\subsection{What Integration Cost, Measured}

Honesty about the bill: the substrate gave vectors durability, transactions,
and replication for free, but not specialist performance for free. Three
costs surfaced in our own measurements (\S\ref{sec:eval},
\S\ref{sec:lessons}): graph index construction at 10M vectors requires several
times the raw vector size in heap because build-time structures hold vectors
on-heap alongside JVector's construction state; sparse query latency on real
SPLADE term distributions degrades at corpus sizes where specialist engines
do not; and settle steps (compaction) that specialists hide behind background
services are visible operations here. We report these as findings with the
same prominence as the wins, and several became upstream fixes during the
course of this work (\S\ref{sec:lessons}).

\subsection{A Note on Measurement Integrity}

The sparse index's original design documentation reported a 766$\times$
speedup over brute-force scoring at 1M documents, a number that survived
review because the same document also \emph{retracted} an earlier
41{,}613$\times$ claim traced to a broken baseline. We mention this not as
trivia but as method: every headline number in this paper carries its
measurement protocol, and where our own instruments were later found
misleading (a memory metric that counted reclaimable page cache; a timeout
that was our harness's fault) the corrected number replaced the flattering
one. \S\ref{sec:eval} states the protocol once, and the raw results are
published with the paper.

\section{Query and Deployment Surfaces}\label{sec:surfaces}
\todo{One executor + result model behind SQL/Cypher/GraphQL/Mongo (Gremlin via
TinkerPop over the same storage); embedded in-process vs server; six wire
protocols; cite the SciPy paper for Python binding internals.}

\section{High Availability}\label{sec:ha}
\todo{Ratis-based Raft shipping shared-WAL page diffs => model-agnostic
replication; quorum/membership hardening lessons.}

\section{Evaluation}\label{sec:eval}
\todo{E1 sparse DB-vs-DB (Qdrant, Milvus) at 100k/1M/10M; E2 unified-ACID hybrid
transaction vs composed stack incl. atomicity demonstration; E3 durability
(kill -9) + Raft failover; E4 embedded vs server deployment axis; E5 dense
sanity. Protocol: docker, pinned digests, identical cpuset/memory caps, N=5
mean$\pm$std, serial re-runs for all reported numbers.}

\section{Lessons Learned and Tradeoffs}\label{sec:lessons}
\todo{The cost of unification (measured); where specialists win; single-node
ceiling; async-ANN consequences; production war stories (co-authors).}

\section{Related Work}\label{sec:related}
\todo{Multi-model (Cosmos DB API-level vs our storage-level; ArangoDB;
Lu \& Holubov\'a survey; UniBench/M2Bench), vector DBs (Milvus, Manu, Qdrant,
pgvector), sparse retrieval (SPLADE, BlockMax-WAND), embedded engines
(SQLite, DuckDB).}

\section{Conclusion}\label{sec:conclusion}
\todo{}

\section*{Acknowledgment}
We thank Christian Himpe (University of M\"unster) for his contributions to
ArcadeDB, for sharing his experience operating it in
DatAasee~\cite{himpe2024dataasee}, and for feedback on this manuscript.
\todo{AI-generated-content disclosure REQUIRED by ICDE 2027: name tools and
affected content areas (drafting/editing assistance, experiment tooling).
Keep updated as the paper evolves.}

\bibliographystyle{IEEEtran}
\bibliography{references}

\end{document}
